\section{Data Sources}
\label{sec:data-sources}

This report describes the final scope of AHI. In that scope the platform draws on three families of data. The first and main family is \textbf{operational data from equipment in use on the International Space Station (ISS)}: wearables, medical equipment, exercise equipment and environmental monitors. The second is \textbf{space weather notifications} from NASA, which warn of solar events before or while their radiation reaches the spacecraft. The third is \textbf{public NASA research data}, used to ground the analysis in real measurements and to support the development and evaluation of the models. \ref{tab:data-sources} summarizes both families and the role of each.

\medskip\noindent\textbf{Design premise.} AHI assumes that every operational value comes from the ISS equipment described in \ref{sec:hardware-sensors}, for both the crew and the cabin. Only equipment that public NASA, Canadian Space Agency (CSA) or Japan Aerospace Exploration Agency (JAXA) sources describe as in use on the ISS is part of this premise. The premise defines what the platform is designed to receive. It does not mean that the project has access to these devices or to their data. Access to real crew data would require authorization, documentation of each interface and a data handling agreement, and it is treated in \ref{sec:limitations-ethics}.

\begin{table}[htbp]
\centering
\vspace{0.8em}
\caption{Data families in the final scope of AHI and their role.}
\label{tab:data-sources}
\small
\begin{tabular}{p{3.4cm}p{5.2cm}p{5.4cm}}
\hline
\textbf{Family} & \textbf{Examples} & \textbf{Role in AHI} \\
\hline
Wearables & Bio-Monitor, Actiwatch Spectrum, Dry-EEG Headband & Continuous or session based physiology, activity and sleep \\
Medical equipment & BP/ECG monitor, Holter Monitor 2, pulse oximeter, Ultrasound 2, optical coherence tomography, fundoscope, vision screening, On-Orbit Hearing Assessment, urine analyzer, blood and saliva collection & Scheduled and periodic clinical measurements, eye and hearing examinations \\
Exercise equipment & CEVIS, ARED, T2 with COLBERT & Physical performance, training load and adherence \\
Environmental monitors & Air Quality Monitor, CO\textsubscript{2} monitors, combustion product analyzer, acoustic dosimeters and sound level meters, radiation monitors and dosimeters, water and microbial monitoring & Context for interpreting physiology; cabin hazards; personal noise and radiation exposure \\
Space weather notifications & DONKI notifications of solar flares, coronal mass ejections and solar energetic particle events, through the NASA Open APIs & Early warning for the radiation shelter recommendation \\
Self report on station computers & Daily check-in, sleep report, cognition tests & Behavioral health, fatigue and cognition \\
Public NASA research data & OSDR, EDA, RadLab, NLSP, LSDA & Reference ranges, model development and retrospective evaluation \\
\hline
\end{tabular}
\vspace{1.2em}
\end{table}

\subsection{Operational Data from ISS Equipment}
\label{subsec:operational-data}

The medical structure of the ISS is the Crew Health Care System. It brings together health maintenance, environmental monitoring and exercise equipment, and the JAXA medical operations page describes its main components \cite{jaxa_medops_system}. AHI organizes the operational data into four categories that follow this structure.

\begin{itemize}
\item \textbf{Wearables} collect physiological data during daily activities without stopping the work of the crew. They are the main source of continuous data.
\item \textbf{Medical equipment} produces clinical measurements at scheduled times, such as blood pressure, an electrocardiogram (ECG) recording, an ultrasound examination or a urine analysis.
\item \textbf{Exercise equipment} describes the daily countermeasure sessions that crews perform to limit the loss of bone, muscle and cardiovascular fitness \cite{nasa_jsc_exercise_countermeasures}.
\item \textbf{Environmental monitors} describe the cabin. They do not measure health directly, but air quality, carbon dioxide, noise and radiation affect the crew and are needed to interpret physiological changes \cite{nasa_toxicology_envchem,nasa_ochmo_tb035_acoustics}. Acoustic dosimeters and personal radiation dosimeters are worn by the crew, so they also measure the exposure of each person.
\end{itemize}

Not every device is used by every crew member during the whole mission. Several wearables are research hardware used during specific investigations \cite{nasa_wearable_tech_iss,nasa_sleep_on_station}. AHI therefore treats each data stream as present only when the corresponding device is in use, and a stream that is expected but absent is reported as missing data and never as normal.

\medskip\noindent\textbf{Self report.} Mood, perceived stress, workload and sleep reports are not produced by hardware. NASA Standard Measures already collects sleep reports and cognition tests in flight \cite{nasa_standard_measures}. AHI assumes that such inputs are entered on station computers or on the tablet of the Bio-Monitor system and are labeled as self report.

\subsection{Collection Frequency}
\label{subsec:collection-frequency}

Not every indicator is collected at the same rate. \ref{tab:frequency} lists the cadence the platform is designed to handle for each kind of data. These cadences are integration targets for AHI. They are not a NASA operational specification.

\begin{table}[htbp]
\centering
\vspace{0.8em}
\caption{Cadence that AHI is designed to handle for each kind of data. Design targets, not a NASA specification.}
\label{tab:frequency}
\small
\begin{tabular}{p{4.6cm}p{4.2cm}p{5.2cm}}
\hline
\textbf{Kind of data} & \textbf{Cadence in AHI} & \textbf{Typical source} \\
\hline
Heart rate & Continuous or near continuous & Bio-Monitor \\
Breathing & Continuous or per session & Bio-Monitor \\
Blood oxygen saturation & Periodic or continuous, depending on the device & Bio-Monitor, pulse oximeter \\
ECG & Per session or 24 hour recording & BP/ECG monitor, Holter Monitor 2 \\
Sleep & Each rest period & Actiwatch Spectrum, Dry-EEG Headband \\
Physical activity & During the day & Actiwatch Spectrum, Bio-Monitor \\
Exercise performance & Per session & CEVIS, ARED, T2 \\
Blood pressure & Scheduled measurements & BP/ECG monitor, Bio-Monitor \\
Cabin CO\textsubscript{2} and air quality & Continuous & CO\textsubscript{2} monitors, Air Quality Monitor \\
Radiation & Continuous area dose; accumulated personal dose & Radiation monitors, personal dosimeters \\
Blood, saliva and urine & Scheduled collections & Collection kits, urine analyzer \\
Ultrasound & Scheduled examinations & Ultrasound 2 \\
Eye examination & Before flight, around flight days 30 and 90, 30 days before return, after flight & Vision screening, fundoscope, optical coherence tomography, Ultrasound 2 \\
Visual symptoms & Any day, in the daily check-in & Self report \\
Hearing thresholds & Periodic, for flights longer than 30 days & On-Orbit Hearing Assessment \\
Personal noise exposure & 24 hour crew-worn sessions; area surveys & Acoustic dosimeters, sound level meter \\
Space weather & On each new notification & DONKI, relayed by the ground node \\
Microbial and water quality & Scheduled sampling & Environmental Health System \\
\hline
\end{tabular}
\vspace{1.2em}
\end{table}

A consequence of these cadences is that each signal has its own definition of fresh and stale. A heart rate that has not arrived for several minutes is stale, while a blood sample result is expected only at the next scheduled collection. AHI stores the expected cadence of every stream and uses it to decide when data become missing.

\subsection{Space Weather Notifications}
\label{subsec:space-weather}

Personal dosimeters measure the dose already received. A solar particle event can raise the dose rate within hours, and NASA-STD-3001 requires real-time monitoring of space weather and a method to alert the crew, locally and remotely, when radiation levels are expected to exceed acceptable levels \cite{nasa_ochmo_tb020_radiation}. AHI therefore adds a second, external source of radiation information: the \textbf{Space Weather Database Of Notifications, Knowledge, Information (DONKI)} of the NASA Community Coordinated Modeling Center \cite{nasa_donki}. DONKI catalogs space weather events, such as solar flares, coronal mass ejections, solar energetic particle events and geomagnetic storms, together with the notifications issued by its forecasters.

\medskip\noindent\textbf{Access through the NASA Open APIs.} DONKI is available through the NASA Open APIs at \texttt{api.nasa.gov} \cite{nasa_open_apis}, with one endpoint per event type and one for notifications. Every request carries an API key. NASA issues personal keys free of charge, and a shared \texttt{DEMO\_KEY} exists for exploration with a much lower rate limit. When tested by the team on 2026-11-15, the notifications endpoint returned current messages, and the event endpoints rejected date ranges longer than 60 days, so the adapter requests short windows. The key is a secret of the ground node: it is stored in the server configuration, never in the interface code, in the repository or in this report, and it can be rotated without a new release.

\medskip\noindent\textbf{Useful, not essential.} DONKI is an internet service. AHI uses it when a connection exists, but its essential functions must keep working without one (\ref{sec:software-architecture}). Only the ground node calls the API. It filters the notifications that matter for the crew, stores them with their message identifier and issue time, and sends them on board through the same store and forward link as any other message. On board, the radiation rule does not wait for DONKI: a sudden rise in the dose rate measured by the station radiation monitors triggers the same recommendation by itself (\ref{sec:analytical-methodology}). DONKI gives earlier warning when the link exists, and the on board monitors give protection when it does not.

\medskip\noindent\textbf{An honest label.} DONKI messages state that the NOAA Space Weather Prediction Center is the official source of space weather forecasts for the United States Government. In a real mission, AHI would receive its alerts from the space weather support of the mission. DONKI is the public NASA source that the design and the demonstration use, and every notification keeps its source and message identifier.

\subsection{Public NASA Research Data}
\label{subsec:public-data}

No public source offers continuous physiological telemetry of astronauts. Public NASA repositories nevertheless hold real measurements that AHI uses in three ways: as historical reference shown next to crew data, as material to develop and test the analysis, and as a retrospective test of the engine outside simulated data.

\begin{itemize}
\item \textbf{NASA Open Science Data Repository (OSDR).} The main NASA repository for space life science, with studies, files and metadata \cite{nasa_about_osdr}. Among its human studies are the Inspiration4 complete blood count (OSD-569) and comprehensive metabolic panel (OSD-575), with four crew members at three time points before launch and four after return and no in-flight samples \cite{osdr_osd569,osdr_osd575}. They give AHI real laboratory values with units, the format that blood collection results take in the platform.
\item \textbf{Environmental Data Application (EDA).} ISS cabin carbon dioxide, temperature, humidity and radiation dose for the periods of research payload missions \cite{osdr_eda}. AHI uses it to replay real cabin periods. The data describe the cabin during a payload mission and not a given crew.
\item \textbf{RadLab.} An access point for radiation telemetry from several instruments \cite{osdr_radlab}. It was unavailable when tested on 2026-11-14, so its coverage and format are still to be checked.
\item \textbf{NASA Life Sciences Portal (NLSP)} and \textbf{Life Sciences Data Archive (LSDA).} Discovery portals for historical life science and human health data, including analog studies \cite{nasa_nlsp,nasa_lsda}. Many datasets require a formal request. They are the path to sleep, cognition and analog data. Analog participants are never presented as astronauts.
\item \textbf{NASA Standard Measures} and \textbf{Artemis~II Standard Measures.} Reference material and not datasets. They list the measures collected before, during and after missions and justify the catalog of indicators \cite{nasa_standard_measures,nasa_artemis_ii_standard_measures}.
\end{itemize}

A caution applies to all public data. It is not safe to assume that individual medical data of astronauts are publicly available, and in general they are not. Public datasets are small, often from short missions, and rarely include in-flight samples. They support grounding and testing, but they are not enough to train clinical models by themselves (\ref{sec:artificial-intelligence}).

\subsection{Data Provenance}
\label{subsec:data-provenance}

Every observation carries a \texttt{data\_origin} field, shown in the interface as a badge, so that a reader always knows where a value comes from:

\begin{itemize}
\item \texttt{DEVICE}: received from ISS equipment through a device adapter (\ref{sec:hardware-sensors}).
\item \texttt{MANUAL}: entered by a person, such as the daily check-in or a sleep report.
\item \texttt{LAB}: a result of a sample collected in flight and analyzed later, with the collection time and the result time kept apart.
\item \texttt{DERIVED}: computed by the platform, such as baselines, risk scores and forecasts.
\item \texttt{REAL\_PUBLIC}: public NASA data, such as research datasets shown as historical reference and space weather notifications from DONKI.
\item \texttt{SIMULATED}: produced by the simulator used for development, testing and demonstration. Simulated values are never mixed with crew data.
\end{itemize}

Each record also carries its source device, unit, measurement time, quality status and activity context. Imported public datasets follow a reproducible path: the original file is downloaded from an allow list of NASA hosts, its SHA-256 hash and download time are recorded, and a manifest keeps the source address, version, population, design and license. Values are kept exactly as reported, nothing is imputed, and fields the source does not state are recorded as not stated.
